\subsubsection{Iterative Prompt Tuning: Sentiment Classification}
\label{sec:experiment-sentiment}

To evaluate the framework's prompt optimization capabilities, we applied it to binary sentiment classification on the Stanford Sentiment Treebank (SST-2)~\cite{socher2013recursive} using \texttt{gpt-4.1-nano}, a lightweight model where prompt engineering has measurable impact. The system was configured with a \texttt{prompt\_tune} task type, a 2\% minimum improvement threshold (\texttt{min\_delta}${}= 0.02$), and a maximum of 5 optimization rounds. Following the framework's train/eval split protocol, we sampled a balanced subset of 20 training examples (10 per class) and 12 held-out evaluation examples (6 per class), drawn from SST-2's training and validation splits respectively using a fixed random seed. All prompt modifications were informed exclusively by training-set failures; the evaluation set was reserved for accept/reject decisions, and leakage checks were enforced at each round.

The framework achieved perfect evaluation accuracy ($1.00$) in 3 rounds, improving from a $0.83$ baseline---without any model fine-tuning or architectural changes.

\paragraph{Round 1 (Baseline).}
The system initialized a minimal prompt consisting of a generic system instruction (``You are a sentiment classifier\ldots Respond with only the sentiment label'') and a zero-shot task template (``Classify the sentiment of the following text as `positive' or `negative'\,'').  This baseline achieved $0.80$ train accuracy (16/20) and $0.8333$ eval accuracy (10/12). Inspection of the 4 training failures revealed two error categories: (i)~ambiguous short fragments lacking explicit sentiment cues (e.g., ``a deeper story,'' ``its opening''), and (ii)~sentences where negation or sarcasm inverted surface-level word polarity.

\paragraph{Round 2 (System Prompt Refinement $\to$ Accepted).}
The investigator agent analyzed the training failures and identified that the model lacked domain context: SST-2 texts are fragments of movie reviews, and phrases like ``a deeper story'' carry positive connotations in film criticism that are absent in general usage. The executor agent revised the system prompt to include three targeted additions:
\begin{enumerate}
  \item \emph{Domain framing}: ``These texts are fragments from movie reviews. Even short or ambiguous phrases should be interpreted in the context of film criticism.''
  \item \emph{Fragment handling}: ``A phrase that describes something approvingly or with admiration is positive, even if it contains words that sound negative in isolation.''
  \item \emph{Negation awareness}: ``Sentences that use negative words (like `isn't', `never', `no') may still express a positive opinion if the overall meaning is favorable.''
\end{enumerate}
This yielded $0.85$ train accuracy (+0.05) and $0.9167$ eval accuracy (+0.0834), exceeding the acceptance threshold. The leakage check confirmed no evaluation examples were referenced.

\paragraph{Round 3 (Few-Shot Examples $\to$ Accepted).}
With 3 training failures remaining, the investigator identified that two of the misclassified fragments (``an era of theatrical comedy that, while past,'' and ``its opening'') were too short and context-dependent for the system prompt alone to resolve. The executor introduced 6 few-shot demonstrations sourced exclusively from the training set, covering three categories:
\begin{itemize}
  \item Short positive fragments: ``making the right choice'' $\to$ positive; ``a deeper story'' $\to$ positive.
  \item Short negative fragments: ``a dumb comedy'' $\to$ negative; ``paper-thin characterizations'' $\to$ negative.
  \item Complex sentences with misleading surface cues, including one with nested negation and one with sarcastic praise.
\end{itemize}
This brought eval accuracy to $1.0000$ (+0.0833) and train accuracy to $0.9000$ (+0.05). The two remaining train failures (``an era of theatrical comedy that, while past,'' and ``its opening'') are arguably mislabeled---both are neutral fragments whose SST-2 ground truth labels are debatable. The framework terminated optimization upon achieving perfect evaluation performance.

Table~\ref{tab:sentiment-progression} summarizes the progression across rounds.

\begin{table}[t]
  \centering
  \caption{Accuracy progression across prompt optimization rounds for SST-2 binary sentiment classification using \texttt{gpt-4.1-nano}. Train ($n{=}20$) and eval ($n{=}12$) subsets are disjoint.}
  \label{tab:sentiment-progression}
  \begin{tabular}{clccl}
    \toprule
    Round & Train Acc. & Eval Acc. & $\Delta_{\text{eval}}$ & Technique \\
    \midrule
    1 (Baseline)  & 0.8000 & 0.8333 & ---      & Zero-shot, generic system prompt \\
    2              & 0.8500 & 0.9167 & $+0.0834$ & Domain-specific system prompt \\
    3              & 0.9000 & 1.0000 & $+0.0833$ & $+$ 6 few-shot exemplars from train \\
    \bottomrule
  \end{tabular}
\end{table}

This experiment demonstrates several properties of the framework in the prompt tuning setting. First, the investigator--executor loop mirrors a practitioner's workflow: diagnose errors on a development set, hypothesize a prompt modification, and validate on held-out data---but executes this cycle autonomously. Second, the system selected complementary strategies across rounds (domain framing, then few-shot examples), consistent with the configured strategy hints (\texttt{["few shots", "chain of thought"]}), though it prioritized system prompt refinement when that alone was sufficient. Third, the strict train/eval separation and per-round leakage checks ensured that the reported eval gains reflect genuine generalization rather than overfitting to the evaluation subset. Finally, the framework self-terminated after 3 of 5 allowed rounds upon reaching perfect eval accuracy, avoiding unnecessary optimization cycles.